\begin{proposition}\label{sga3:XXI:3.6.9}
\begin{enumerate}
\item[\SGAInlineRefTarget{sga3:XXI:3.6.9:item:i}(i)] For every \(x\in X\), \(R_+(x)\) is a system of positive
roots. For every system of positive roots \(R_+\),
\(\mathcal{C}(R_+)\) is a chamber in \(V_{\mathbb R}^\ast\). The
\(\mathcal{C}(R_+)\) are the maximal convex subsets of \(X\).
\item[\SGAInlineRefTarget{sga3:XXI:3.6.9:item:ii}(ii)] Let \(\Delta\) be a system of simple roots. One has
\[
 \mathcal{C}(\mathscr{P}(\Delta))
 =\{x\in V_{\mathbb R}^\ast
       \mid (x,\alpha)>0\text{ for every }\alpha\in\Delta\}.
\]
The walls of \(\mathcal{C}(\mathscr{P}(\Delta))\) are the
hyperplanes
\[
 \mathcal{H}_\alpha=\alpha^{-1}(0),
 \qquad \alpha\in\Delta;
\]
its faces are
\[
 C_\alpha
 =\{x\in V_{\mathbb R}^\ast
      \mid (x,\alpha)=0,\ 
      (x,\beta)>0\text{ for }\beta\in\Delta,\ \beta\ne\alpha\}.
\]
\item[\SGAInlineRefTarget{sga3:XXI:3.6:item:iii}(iii)] One has the equivalence
\[
 R_+(x)=R_+
 \quad\Longleftrightarrow\quad
 x\in\mathcal{C}(R_+).
\]
\end{enumerate}
\end{proposition}

\begin{proof}
It is first clear that \(R_+(x)\) is a system of positive roots.
Since \(x\in\mathcal{C}(R_+(x))\), the union of the
\(\mathcal{C}(R_+(x))\) is \(X\). Property \SGARefLink{sga3:XXI:3.6:item:iii}{(iii)} is immediate; it
follows that the \(\mathcal{C}(R_+)\) form a partition of \(X\).
The first assertion of \SGARefLink{sga3:XXI:3.6.9:item:ii}{(ii)} is evident. It follows at once that
\(\mathcal{C}(R_+)\) is a chamber in \(V_{\mathbb R}^\ast\), which
proves the remainder of \SGARefLink{sga3:XXI:3.6.9:item:i}{(i)}. Put
\[
 C=\mathcal{C}(\mathscr{P}(\Delta)).
\]
It remains only to note that
\[
 \overline C\cap
 \left(V_{\mathbb R}^\ast\setminus
       \bigcup_{\alpha\in R}\alpha^{-1}(0)\right)
 =
 \overline C\cap
 \left(V_{\mathbb R}^\ast\setminus
       \bigcup_{\alpha\in\Delta}\alpha^{-1}(0)\right)
\]
to complete the proof of \SGARefLink{sga3:XXI:3.6.9:item:ii}{(ii)} by
Lemma~\ref{sga3:XXI:3.6.1}.
\end{proof}

\begin{definition}\label{sga3:XXI:3.6.10}
The \(\mathcal{C}(R_+)\) are called the \emph{Weyl chambers} of the
root datum. Conversely, for every Weyl chamber \(C\), put
\[
 R_+(C)=R_+(x)
\]
for any \(x\in C\).
\end{definition}

\begin{corollary}\label{sga3:XXI:3.6.11}
The maps
\[
 R_+\longmapsto\mathcal{C}(R_+),
 \qquad
 C\longmapsto R_+(C)
\]
give a bijective correspondence between systems of positive roots
and Weyl chambers.
\end{corollary}

This correspondence is invariant under the Weyl group:

\begin{lemma}\label{sga3:XXI:3.6.12}
If \(w\in W(\mathcal{R})\), then
\[
 \mathcal{C}(w(R_+))=w(\mathcal{C}(R_+)).
\]
\end{lemma}

\begin{corollary}\label{sga3:XXI:3.6.13}
The correspondences
\[
 \Delta\longleftrightarrow R_+\longleftrightarrow C
\]
are isomorphisms of homogeneous spaces under \(W(\mathcal{R})\).
\end{corollary}

We shall see later that these homogeneous spaces are principal
(\SGARefLink{sga3:XXI:5.5}{5.5}).

\begin{remark}\label{sga3:XXI:3.6.14}
If \(C\) is a Weyl chamber, then \(-C\) is also one; it is called the
\emph{opposite} of \(C\). Consequently there exists
\(w_0\in W\)---and in fact exactly one, cf.
\SGARefLink{sga3:XXI:5.5}{5.5}---such that
\[
 w_0(C)=-C.
\]
It is called the \emph{symmetry} of the root datum relative to the
Weyl chamber \(C\) (or relative to \(R_+(C)\), or to
\(\mathscr{S}(R_+(C))\), and so on).
\end{remark}
